本书源自笔者在雷根斯堡大学讲授的两门课程的讲义：一门是 2018/2019 冬季学期的
\emph{Introduction to ∞-categories}，另一门是 2019 夏季学期的 \emph{∞-categories}。
书中大部分材料并非原创；在本书前几章的许多地方，我都沿袭了\trRezk{}（Rezk）的论证\cite{Rez20}。
当然，进一步的灵感来源还有\trLurie{}的著作\emph{Higher Topos Theory}\footnote{译注：本书所引书名与课程名一律保留原文。}（\cite{Lur09}）、
\trCisinski{}的著作\emph{Higher Categories and Homotopical Algebra}（\cite{Cis19}）、
\trJoyal{}关于\trQuasiCategory{}的论文（\cite{Joy08}），以及%
\trHaugseng{}（Haugseng）的\trInfinityCategories{}讲义（\cite{Hau17}）。
此外还有两篇综述文章：拉恩（Rahn）的\emph{A short course on ∞-categories}
（本书写作时或署名格罗特（Groth）），以及\trGepner{}（Gepner）的\emph{An introduction to higher categorical algebra}，
二者均作为一章收入\cite{Mil19}。特别是，\trGepner{}的这篇论文十分好读，
适合在攻读\trLurie{}关于这一主题的开创性著作\emph{Higher Algebra}（\cite{Lur17}）之前先行阅读。

阅读本书时，若对普通\trCategoryTheory{}已有一些接触将大有裨益，但这并非严格必需：本书所需的
先修知识并不多，而且最重要的概念都会在卷首予以回顾。例如，见过一些%
\trCategory{}、\trFunctor{}、\trAdjunction{}、\trColimit{}与\trLimit{}的例子会很有用；
对\trAlgebraicTopology{}的某些了解——例如基础\trHomotopyTheory{}，以及\trTopologicalSpace{}与\trSimplicialSets{}之间的关系——也会很有帮助。
特别是，知道\trAlgebraicTopology{}中出现的许多\trCategory{}都容许\trTopologicalSpace{}上的\trEnrichment{}（enrichment），
将有助于体会\trCoherentNerve{}（coherent nerve）的构造——它从每个在\trTopologicalSpace{}中\trEnrichment{}的\trCategory{}%
得出相应的\trInfinityCategory{}；而且，熟悉\trModelCategory{}的读者固然会在本书前两章中
见到一些熟悉的论证，但要跟上本书的推理，并不需要任何\trModelCategory{}的知识。

总而言之，对于在\trHomotopyTheory{}与\trCategoryTheory{}方面已有一些背景、并对\trInfinityCategories{}的
基础理论怀有兴趣的读者而言，本书所呈现的材料是非常合适的。本书可以用作该主题
一门讲授课的教材，涵盖约两个学期、每周 4 小时的讲授；书中这组习题可供配套的
习题课使用。同样，本书也可用于研读课，或者，当然，作为个人读物来自行阅读。

丹麦哥本哈根\hfill\trLand
